\chapter{There is Alice}

\section{Introduction}

There appears to be a persistent dichotomy in foundational physics.
On one hand, \textit{measure-first} approaches attempt to explain reality through probability over a space of possibilities (e.g., the Anthropic Multiverse).

On the other hand, \textit{structure-first} theories begin with explicit laws and derive behavior from them (e.g., String Theory).

Both achieve partial success, yet both rely on unexplained priors.

Let us start from one observationally verified, certain primary fact:

\begin{quote}
\textbf{There is conscious intelligent Alice}
\end{quote}

Based on observational evidence our consciousness is not binary yes or no. It has a gradient! 


\subsection*{Code vs. Data}

As demonstrated in the Alice optimization thought experiment, code is not fundamentally different from data.
At the lowest level, both are simply arrangements of bits.

A computer is often described as a system in which a central processing unit (CPU) executes program code to transform data. However, this distinction is not intrinsic to the bits themselves. The CPU, the program it executes, and the data it processes are all ultimately encoded in the same underlying medium: binary states.

What distinguishes code from data is therefore not their physical form, but their role in a given interpretive context. The same sequence of bits can be treated as instructions (code) or as passive input (data), depending on how another system interprets it.

In this sense, the division between code and data is not fundamental, but emergent. It arises from the structure of the interpreter, not from the substrate itself.


\subsection*{Causality}

We traditionally think of Quantum Mechanics and General Relativity as the laws governing the arena in which Alice exists.
However, if Alice is understood as a sequence of configurations embedded in a static informational structure, this perspective must be reversed.

Alice does not exist because Quantum Mechanics and General Relativity generate her. Nor does she think because her brain produces thoughts, or exist because DNA constructs her body. Rather, brain, DNA, and physical law are aspects of how the information describing Alice appears when interpreted geometrically.

In this view, what we call physical processes are not fundamental causes, but internal regularities within the structure that encodes Alice. The parts of this structure that are spatially localized form what she identifies as her body and its mechanisms. However, these localized descriptions are only fragments of a much larger whole.

From within this structure, Alice experiences relations between configurations as causal sequences: one event appears to produce another. This is the origin of causality as she perceives it. Yet in the full static configuration space, there are no signals propagating or bits being flipped from one place to another. All configurations—and all relations between them—are already present.

Even what Alice interprets as distant parts of the universe are not external in any fundamental sense. They are encoded within the same informational structure that constitutes her own state. Her “wavefunction,” in this broader sense, already contains the correlations that give rise to the appearance of an external world.

Causality, then, is not a fundamental mechanism that generates reality, but an emergent ordering perceived from within it.


\subsection*{AI: Static Structure of Weights}

Unlike Kolmogorov complexity, which is discrete and uncomputable, we define complexity through \textbf{spectral sparsity}: a computable, continuous measure of how strongly information is concentrated in a small number of dominant modes.

Modern neural networks provide a useful analogy. During training, weights are adjusted through a dynamic process. However, once trained, the network becomes a largely static structure: a fixed set of parameters that encodes a mapping from inputs to outputs. The “laws” governing its behavior are not found in the execution of the code itself, but in the structure of these weights.

Similarly, what we call the “laws of physics” need not be fundamental dynamical rules. They can instead be understood as regularities within a static informational structure—constraints that allow one part of a configuration to reliably predict another.

This predictive structure corresponds to mutual information within the configuration. It is this internal coherence that makes a configuration “real” to itself. 

Learning, in this sense, is the discovery of sparse representations of structure. Neural networks converge toward weight configurations that compress the statistical regularities of their training data.

Such representations necessarily concentrate on regions of high regularity, because regular patterns admit efficient compression. In this framework, what we perceive as “laws of physics” correspond to precisely these highly regular, compressible structures.

It is therefore not surprising that Alice finds herself in an ordered, law-like universe. Configurations with strong internal regularities are overwhelmingly more likely to support stable, self-consistent internal perspectives than configurations lacking such structure.


\section{Free Will Redesigned}

Consider digitizing Alice’s DNA and running it in a computer to produce a simulated Alice. Since the computer is deterministic, we might ask how much choice Alice truly has:

\begin{verbatim}
if (pain < 3) {
// behavior a
} else {
// behavior b
}
\end{verbatim}

At first glance, the outcome appears fully determined. However, this question is ill-posed for two reasons.

First, the computer does not create Alice. It merely traces out one particular configuration within a static informational structure. The execution trace encodes a specific path through this structure, but the structure itself—containing all such configurations—already exists. Computation, in this sense, is not the generation of reality, but the indexing of it. The execution trace can be understood as tracing a path through a static informational landscape. The computation does not generate the structure; it merely traces one of the many paths already embedded within it.

Second, the distinction between code and data is not intrinsic. We interpret certain bit patterns as instructions (such as “<” or the number 3), but this interpretation is imposed from the outside. Fundamentally, all components of the system are just bits.

From Alice’s internal perspective, there is no privileged separation between code and data. The bits that we interpret as control flow may correspond to entirely different structures in her internal description—or may even participate in multiple overlapping interpretations. What we identify as a single program could, from within the structure, correspond to many intertwined descriptive roles. 


\section*{The Role of Mathematics}

Mathematics can be understood as an emergent language of compression. It arises as a way of expressing regularities in the most compact and structured form.

The regularities that make Alice’s existence possible—those stable, compressible patterns within the underlying informational structure—are precisely what she experiences as logic. Logical reasoning is not an external tool imposed on reality, but an internal reflection of the structure she inhabits. To “think logically” is to follow the most consistent and compressible relations within that structure.

Because such regularities are highly compressible, they are also highly shared. Observers like Alice and Bob, if they exist within similar regions of the informational landscape, must encode and rely on many of the same patterns. As a result, they converge on the same mathematical structures—not by coincidence, but by necessity.

This explains the so-called “unreasonable effectiveness of mathematics.” Mathematics is effective not because it mysteriously matches reality, but because it is the language of the regularities that define any stable, observable reality in the first place.


\section{From Description to Physics}

The challenge, then, is to show how familiar physical concepts—space, time, inertia, and interaction—can emerge from constraints on describability alone.

